%! AP Statistics — Unit 4, Lesson 4.4 — Guided Notes
\documentclass[10pt]{article}
\usepackage{apstats-article}
\usepackage{apstats-boxes}

\begin{document}

\pageheader{Unit 4, Lesson 4.4}{Guided Notes}
\namedateperiod

\vspace{0.12in}

\begin{objectivebox}
\small By the end of this lesson, I will be able to\dots\\[4pt]
\writeline\\[1pt]
\writeline\\[1pt]
\writeline
\end{objectivebox}

\vspace{0.1in}

\begin{vocabbox}
\termblanklong{Joint probability $P(A \cap B)$}
\termblanklong{Intersection $A \cap B$}
\termblanklong{Mutually exclusive (disjoint) events}
\termblanklong{Sample space}
\end{vocabbox}

\vspace{0.1in}

\begin{hookbox}
\small
\textbf{Two scenarios:}
\begin{enumerate}[leftmargin=1.5em, itemsep=4pt]
  \item A randomly selected student is enrolled in both AP Statistics and AP Calculus.
  \item A randomly selected card from a standard deck is both a heart and a spade.
\end{enumerate}
\textbf{Which is impossible? Why?}\\[4pt]
\writeline\\[1pt]\writeline
\end{hookbox}

\vspace{0.1in}

\begin{notesbox}{Joint Probability and Intersections}
\small
\textbf{Definition:} The \textbf{joint probability} $P(A \cap B)$ is the probability that
event $A$ \blank{4cm} and event $B$ \blank{4cm} simultaneously.

\vspace{10pt}
\textbf{Notation:} $A \cap B$ is read as ``$A$ \blank{2.5cm} $B$'' and represents all
outcomes in \blank{4cm}.

\vspace{10pt}
\textbf{Two Venn Diagrams:}

\vspace{6pt}
\begin{minipage}[t]{0.46\linewidth}
  \centering
  \begin{tcolorbox}[colback=white, colframe=navy, arc=2mm, width=\linewidth,
      left=2mm, right=2mm, top=2mm, bottom=2mm]
    \centering\small
    \textbf{Not Mutually Exclusive}\\[4pt]
    \begin{tikzpicture}
      \draw[navy,thick] (0,0) circle (0.7cm) node[left=0.55cm,navy]{\small$A$};
      \draw[navy,thick] (0.9,0) circle (0.7cm) node[right=0.55cm,navy]{\small$B$};
      \fill[sky,opacity=0.5] (0.45,0) ellipse (0.26cm and 0.55cm);
      \node at (0.45,0) {\tiny$A\cap B$};
    \end{tikzpicture}\\[4pt]
    $P(A \cap B)$ \blank{1.5cm}
  \end{tcolorbox}
\end{minipage}
\hfill
\begin{minipage}[t]{0.46\linewidth}
  \centering
  \begin{tcolorbox}[colback=white, colframe=navy, arc=2mm, width=\linewidth,
      left=2mm, right=2mm, top=2mm, bottom=2mm]
    \centering\small
    \textbf{Mutually Exclusive}\\[4pt]
    \begin{tikzpicture}
      \draw[navy,thick] (-0.85,0) circle (0.7cm) node[navy]{\small$A$};
      \draw[navy,thick] (0.85,0) circle (0.7cm) node[navy]{\small$B$};
    \end{tikzpicture}\\[4pt]
    $P(A \cap B) = $ \blank{1.5cm}
  \end{tcolorbox}
\end{minipage}
\end{notesbox}

\vspace{0.1in}

\begin{notesbox}{Mutually Exclusive Events}
\small
\textbf{Definition (VAR-4.C.2):} Two events $A$ and $B$ are \textbf{mutually exclusive}
(also called \textbf{disjoint}) if \blank{6cm}.\\[6pt]
Therefore, $P(A \cap B) = $ \blank{2cm}.

\vspace{10pt}
\textbf{How to check:}
\begin{enumerate}[leftmargin=1.5em, itemsep=4pt]
  \item List the outcomes in event $A$: \blank{6cm}
  \item List the outcomes in event $B$: \blank{6cm}
  \item Ask: does any single outcome appear in \blank{4cm}?\\
        If YES $\to$ \blank{3.5cm} \quad If NO $\to$ \blank{3.5cm}
\end{enumerate}

\vspace{10pt}
\textbf{Important:} Mutually exclusive $\neq$ \blank{4cm} (coming in Lesson 4.6!).
\end{notesbox}

\vspace{0.1in}

\begin{notesbox}{Worked Examples}
\small
\textbf{Example 1:} Roll a fair die. $A = \{\text{odd}\}$, $B = \{\text{even}\}$.
\begin{itemize}[itemsep=3pt]
  \item Outcomes in $A$: \blank{4cm}
  \item Outcomes in $B$: \blank{4cm}
  \item $A \cap B = $ \blank{2cm} \quad So: mutually exclusive? \blank{2cm}
  \item $P(A \cap B) = $ \blank{2cm}
\end{itemize}

\vspace{8pt}
\textbf{Example 2:} Draw a card. $A = \{\text{heart}\}$, $B = \{\text{face card}\}$.
\begin{itemize}[itemsep=3pt]
  \item Outcomes in both: \blank{6cm} (name them)
  \item $A \cap B = \varnothing$? \blank{2cm} \quad So: mutually exclusive? \blank{2cm}
  \item $P(A \cap B) = $ \blank{2cm}
\end{itemize}

\vspace{8pt}
\textbf{AP Exam tip:} Always write a complete sentence justification:\\[4pt]
``Events $A$ and $B$ \blank{3cm} mutually exclusive because \blank{5cm}.''
\end{notesbox}

\vspace{0.1in}

\begin{practicebox}
\small For each pair, (1) list any shared outcomes, (2) state whether the events are
mutually exclusive, (3) justify your answer in context.

\vspace{6pt}
\begin{enumerate}[leftmargin=1.5em, itemsep=10pt]
  \item A student is randomly selected. $A = \{\text{sophomore}\}$, $B = \{\text{on the honor roll}\}$.\\
        Shared outcomes: \blank{4cm} \quad ME? \blank{2cm}\\
        Justification: \writeline

  \item A month is randomly chosen. $A = \{\text{has 31 days}\}$, $B = \{\text{is in summer}\}$.\\
        Shared outcomes: \blank{4cm} \quad ME? \blank{2cm}\\
        Justification: \writeline

  \item An integer from 1--10 is chosen. $A = \{\text{multiple of 3}\}$, $B = \{\text{multiple of 5}\}$.\\
        Shared outcomes: \blank{4cm} \quad ME? \blank{2cm}\\
        Justification: \writeline
\end{enumerate}
\end{practicebox}

\end{document}
